% !TEX TS-program = xelatex
% !TEX encoding = UTF-8 Unicode
% ============================================================================
%  GUION DE LECTURA — Seminario de Investigación I (Avance)
%  Acompaña a talk_seminario_avance.tex (mismos títulos de frame).
%  Texto hablado, fluido, para leer en voz alta durante la presentación online.
%
%  El número al inicio de cada título coincide con el número de diapositiva
%  que aparece en el pie del talk principal, para sincronizar fácilmente.
%
%  Compila con XeLaTeX o LuaLaTeX (o tectonic). Formato deliberadamente simple.
% ============================================================================
\documentclass[11pt,aspectratio=169]{beamer}

\usepackage[spanish,es-noshorthands]{babel}
\usepackage{amsmath,amssymb}

% ---------- Formato mínimo y legible ----------
\usetheme{default}
\setbeamertemplate{navigation symbols}{}
\setbeamercolor{frametitle}{fg=black}
\setbeamerfont{frametitle}{series=\bfseries,size=\large}
\setbeamertemplate{frametitle continuation}{\quad(cont.)}

% Texto a la izquierda, interlineado cómodo para leer
\setbeamertemplate{itemize item}{--}
\setlength{\parskip}{0.55em}
\linespread{1.08}

\title{Guion de lectura — Avance}
\subtitle{Efectos de la IA Generativa en los Mercados Laborales de América Latina}
\author{Eric Torres}
\date{Seminario de Investigación I}

\begin{document}

% ===== Portada del guion =====
\begin{frame}[plain]
  \titlepage
  \vfill
  {\footnotesize Documento de apoyo para leer en voz alta. El número de cada
  título corresponde a la diapositiva del talk principal.}
\end{frame}

% ===== Portada del talk =====
\begin{frame}[allowframebreaks]{Portada — Apertura}
Muy buenas tardes a todas y a todos. Muchas gracias por su tiempo y por estar
aquí. Mi nombre es Eric Torres, soy estudiante del Programa Doctoral en Economía
de la Pontificia Universidad Católica del Perú.

Hoy quiero presentarles un avance de mi investigación sobre los efectos de la
inteligencia artificial generativa —en concreto, de los grandes modelos de
lenguaje, los LLMs, que es la tecnología detrás de herramientas como ChatGPT—
sobre los mercados laborales de América Latina.

La pregunta de fondo es muy sencilla de enunciar y muy difícil de responder:
¿esta nueva tecnología ya está cambiando el empleo, los salarios y las horas de
trabajo de la gente? Y, sobre todo, ¿está cambiándolos de la misma manera en una
región como la nuestra, donde el mercado laboral funciona de un modo muy distinto
al de los países ricos?

Quiero ser transparente desde el inicio: esto es un avance. La meta del proyecto
es comparar seis países de la región, pero hoy traigo resultados terminados solo
para el Perú, concretamente para Lima Metropolitana. Sobre esos resultados peruanos
voy a concentrar la mayor parte de la charla.
\end{frame}

% ===== 1 =====
\begin{frame}[allowframebreaks]{1 · Hoja de ruta del avance}
Permítanme primero darles el mapa de la charla, para que sepan a dónde vamos.

Empiezo con la \textbf{motivación}: por qué este es un buen momento para estudiar
esto, y por qué casi toda la evidencia que tenemos viene de fuera de nuestra región.

Luego planteo la \textbf{pregunta} central: ¿movieron los LLMs los resultados del
mercado laboral en América Latina?

Después explico la \textbf{estrategia}: uso una metodología de diferencias en
diferencias con tratamiento continuo, propuesta por Callaway, Goodman-Bacon y
Sant'Anna en 2024, y la aplico país por país.

En la parte central muestro \textbf{lo que traigo hoy}: las estimaciones para el
Perú, usando la Encuesta Permanente de Empleo Nacional, entre el primer trimestre
de 2021 y el cuarto de 2025.

Y cierro con lo que queda \textbf{pendiente}: los otros cinco países —Costa Rica,
Colombia, Ecuador, México y Uruguay—.

Insisto en una idea que repetiré varias veces: esta presentación es un avance.
Solo reporto estimaciones preliminares para Perú. La comparación entre los seis
países queda pendiente, y esa comparación es, en realidad, la contribución
empírica central del artículo.
\end{frame}

% ===== 2 =====
\begin{frame}[allowframebreaks]{2 · Motivación: General-Purpose Technology}
Antes de hablar del mercado laboral, conviene entender qué tipo de tecnología
tenemos entre manos.

Los economistas usamos el término ``tecnología de propósito general'', en inglés
\emph{General-Purpose Technology} o GPT. Una tecnología de propósito general es
aquella que es tan importante que llega a transformar gran parte de la economía y
de la sociedad. No es una herramienta diseñada para una sola tarea, sino más bien
una plataforma de base sobre la que se construyen muchísimas otras cosas.

¿Cómo reconocemos una tecnología de propósito general? Tiene tres características.
Primero, se usa en muchos sectores a la vez, no en uno solo. Segundo, mejora de
forma continua con el tiempo, es decir, evoluciona. Y tercero, genera innovación
complementaria: aparecen nuevos productos y nuevas formas de trabajo alrededor de
ella.

Los ejemplos históricos son los grandes hitos de la economía moderna: la máquina
de vapor, la electricidad, las computadoras y el internet. La tesis de fondo de
este proyecto es que la inteligencia artificial —y muy en particular los modelos
de lenguaje de hoy— pertenece a esa misma categoría. Si eso es cierto, entonces
sus efectos sobre el trabajo no son un tema de moda pasajera: son algo que vale
la pena medir con cuidado.
\end{frame}

% ===== 3 =====
\begin{frame}[allowframebreaks]{3 · Motivación: un shock global, evidencia local}
Concretemos. El 30 de noviembre de 2022 se lanzó ChatGPT. En cuestión de semanas,
puso un modelo de lenguaje de propósito general en manos de cientos de millones de
personas. Fue, literalmente, un shock global y casi instantáneo.

¿Qué tan grande puede ser su impacto potencial? Eloundou y coautores, en un
trabajo de 2024, estiman que alrededor del cincuenta por ciento de los empleos en
Estados Unidos tienen al menos una tarea que podría reducirse en al menos un
cincuenta por ciento usando un LLM. Es una cifra enorme; sugiere que el potencial
de exposición es muy amplio.

Pero —y este es el punto importante— una cosa es el potencial y otra muy distinta
es el efecto realizado. Cuando uno mira las primeras estimaciones causales
rigurosas, los efectos medidos son sorprendentemente pequeños. Humlum y
Vestergaard, en Dinamarca, descartan efectos sobre salarios u horas mayores al
dos por ciento. Hartley y coautores, y Chen y coautores en Estados Unidos,
encuentran efectos detectables, pero concentrados en las ocupaciones más
expuestas. Hui y coautores sí hallan caídas grandes, del veinte al cincuenta por
ciento, pero en plataformas de trabajo \emph{freelance}.

Y aquí está el vacío que motiva mi tesis: casi toda esta evidencia causal proviene
de Estados Unidos, de Dinamarca o de plataformas dominadas por trabajadores de
economías avanzadas. Prácticamente no sabemos nada, de forma causal, sobre lo que
pasa en una economía en desarrollo.
\end{frame}

% ===== 4 =====
\begin{frame}[allowframebreaks]{4 · El shock llegó a la región}
Antes de preguntarnos por los efectos, vale la pena confirmar algo básico: ¿el
shock siquiera llegó a nuestra región?

En esta diapositiva, el gráfico de la izquierda muestra el interés de búsqueda en
Google Trends para los seis países del estudio: Costa Rica, Colombia, Ecuador,
México, Perú y Uruguay. La historia que cuenta es muy clara. En las ocho semanas
previas al lanzamiento de ChatGPT, el interés era prácticamente cero. A lo largo
de 2023 se produce una subida abrupta. Y el máximo se alcanza a inicios de 2025.
Es decir: la curiosidad y el uso explotaron en muy poco tiempo, igual que en el
resto del mundo.

Ahora bien, quiero ser cuidadoso con lo que esto significa, porque es fácil
sobreinterpretarlo. Que la gente conozca y busque ChatGPT no es lo mismo que que
lo adopte en su trabajo, y la adopción tampoco es lo mismo que un efecto causal
sobre el mercado laboral. Reconocimiento, adopción y efecto causal son tres cosas
distintas. Google Trends solo me confirma la primera. El resto de la charla se
dedica, justamente, a intentar medir la tercera.
\end{frame}

% ===== 5 =====
\begin{frame}[allowframebreaks]{5 · ¿Por qué América Latina?}
¿Por qué estudiar América Latina, y no simplemente replicar los estudios de
Estados Unidos? Porque nuestra región tiene rasgos estructurales que pueden hacer
que la tecnología actúe de una manera completamente distinta.

Primero, la \textbf{informalidad}. La tasa promedio de la región ronda el
cincuenta y cinco por ciento, pero con un rango enorme: desde alrededor del
veinticinco por ciento en Uruguay hasta más del setenta por ciento en la región
andina. Esto no existe en los países donde se ha hecho la mayoría de los estudios.

Segundo, son \textbf{mercados segmentados}. El sector formal y el informal
funcionan casi como dos mercados distintos: tienen instituciones diferentes para
fijar salarios, distinta exposición digital y distintos colchones para
amortiguar un shock de productividad.

Tercero, la \textbf{exposición a la tecnología es estructuralmente distinta}. En
un trabajo del BID de 2024 —del que soy coautor, Azuara, Ripani y Torres—
mostramos que la exposición a los LLMs en la región se inclina hacia los
trabajadores formales, urbanos y con educación terciaria.

Y aquí está, de nuevo, el vacío: ningún artículo ha dado el paso de estimar de
forma \emph{causal} si esa exposición efectivamente movió los resultados del
mercado laboral. La literatura latinoamericana que existe —Ciaschi y coautores,
Egana-delSol y Bravo-Ortega— es valiosa, pero es descriptiva sobre exposición.
Describe quién está expuesto, no qué le pasó. Este artículo busca dar ese paso
causal.
\end{frame}

% ===== 6 =====
\begin{frame}[allowframebreaks]{6 · Pregunta de investigación y contribuciones}
Con todo eso, puedo enunciar la pregunta de investigación de forma precisa.

¿Cuál es el efecto causal de la exposición ocupacional a los LLMs sobre el empleo,
los salarios y las horas de trabajo en seis países de América Latina? Y,
crucialmente, ¿cómo varía ese efecto según la formalidad de base de la ocupación?

El proyecto hace cinco contribuciones. Primera: ofrece las primeras estimaciones
causales fuera de la OCDE sobre el efecto laboral de los LLMs. Segunda: aplica el
estimador de diferencias en diferencias con tratamiento continuo de Callaway y
coautores, país por país. Tercera: pone la formalidad de base en el centro del
análisis de heterogeneidad, no como un simple subgrupo de robustez, sino como el
margen principal de interés. Cuarta: construye un \emph{crosswalk} y una base de
datos armonizada para seis países, entre 2021 y 2025, que es reproducible. Y
quinta —esta es metodológica— precompromete cuatro narrativas posibles antes de
mirar los datos, de modo que el patrón que encontremos en cada país lo determine
el dato, y no una búsqueda \emph{ex post} de la especificación que más nos guste.
\end{frame}

% ===== 7 =====
\begin{frame}[allowframebreaks]{7 · Marco teórico: tareas, desplazamiento, reposición}
Déjenme dar el marco teórico, que es sencillo pero ordena todo el análisis. Me
apoyo en el enfoque de producción basada en tareas, de Acemoglu, Autor y Restrepo.

La idea es esta: una ocupación no es una caja negra; es un paquete de tareas. Cada
ocupación dedica una proporción de su tiempo a distintas tareas. La exposición de
una ocupación a los LLMs, que llamo $\beta_o$, es simplemente la suma de las
proporciones de tiempo dedicadas a tareas que un LLM puede complementar. Una
ocupación con $\beta$ alto tiene muchas tareas ``llm-izables''; una con $\beta$
bajo, pocas.

¿Qué le pasa al salario cuando llegan los LLMs? La respuesta neta tiene dos
fuerzas opuestas. Por un lado, un efecto de \textbf{desplazamiento}: la máquina
hace la tarea y reduce la demanda de trabajo. Por otro, un efecto de
\textbf{reposición o productividad}: el trabajador se vuelve más productivo y eso
empuja su salario hacia arriba. El efecto neto sobre el salario es la diferencia
entre reposición y desplazamiento, multiplicada por la exposición $\beta$.

El punto clave es que el \emph{signo} de esa diferencia no es una predicción
teórica: es un objeto empírico. La teoría no me dice si gana la reposición o el
desplazamiento; eso hay que medirlo.

Y aquí entra el canal que más me interesa, el de la formalidad. Adoptar estas
herramientas en una empresa requiere una inversión de costo fijo. Por eso la
adopción tiende a concentrarse en las firmas formales. Como consecuencia, el
efecto en el sector formal y en el informal puede diferir tanto en magnitud como
en signo. Esa intuición es la que después voy a poner a prueba con los datos.
\end{frame}

% ===== 8 =====
\begin{frame}[allowframebreaks]{8 · Cuatro narrativas pre-registradas, país por país}
Esta diapositiva contiene el corazón metodológico del diseño, así que me detengo
un momento.

Antes de mirar los resultados posteriores al tratamiento, dejé por escrito cuatro
historias posibles, cuatro narrativas, con sus predicciones. La idea es
comprometerse de antemano para no caer en la tentación de inventar la explicación
después de ver el número.

La narrativa A, ``la informalidad amplifica'', predice efectos grandes y negativos,
sobre todo en los informales, que absorberían más el golpe. La narrativa B, la
``replicación del nulo'', dice que no pasa casi nada, igual que en Dinamarca: el
Perú se parecería a Humlum y Vestergaard. La narrativa C, ``nulo formal, grande
informal'', predice que los formales no se mueven pero los informales sí sufren un
efecto negativo importante; sería un canal que los estudios de Estados Unidos o
Dinamarca no pueden detectar, porque allá casi no hay informalidad. Y la narrativa
D, ``ganancias estratificadas'', predice ganancias salariales, pero concentradas
en los formales urbanos con educación terciaria.

La contribución metodológica está, precisamente, en precomprometer estas
narrativas antes de mirar el dato post-tratamiento. Ahora bien, seré honesto con
ustedes: las estimaciones de Perú se ejecutaron antes del depósito formal del
pre-registro. Por eso las trato como \textbf{exploratorias} en el mapeo a
narrativas. El pre-registro formal, en firme, cubrirá los cinco países restantes.
\end{frame}

% ===== 9 =====
\begin{frame}[allowframebreaks]{9 · Estrategia empírica: DiD continuo, país por país}
Veamos ahora la estrategia empírica con un poco más de detalle, pero sin perdernos
en lo técnico.

Uso una sola fecha de adopción global: el cuarto trimestre de 2022. Ese trimestre,
el del lanzamiento, lo descarto como ``colchón'', y considero que el tratamiento
empieza de forma efectiva en el primer trimestre de 2023.

La idea de diferencias en diferencias es comparar cómo evolucionan, antes y después
de esa fecha, las ocupaciones más expuestas frente a las menos expuestas. Como la
exposición $\beta$ es una variable continua —no un simple ``tratado o no
tratado''— uso el estimador de tratamiento continuo de Callaway, Goodman-Bacon y
Sant'Anna.

Los resultados que miro son tres, a nivel de celda ocupación-trimestre: el
logaritmo del empleo, el logaritmo del salario por hora, y las horas semanales. La
exposición $\beta$ viene de Eloundou y coautores, y la mapeo a la clasificación
internacional de ocupaciones, la CIUO-08 a cuatro dígitos, con un emparejamiento
del noventa y ocho por ciento. Los errores estándar los agrupo a nivel de
ocupación.

Y, muy importante, no me quedo con una sola especificación. Apoyo el resultado
con varias herramientas complementarias: un \emph{event study} de Sun y Abraham,
una versión binaria del estimador de Callaway y Sant'Anna partiendo en la mediana,
un enfoque tipo \emph{shift-share}, y un \emph{wild cluster bootstrap}. La idea es
que la conclusión no dependa de una única receta.
\end{frame}

% ===== 10 =====
\begin{frame}[allowframebreaks]{10 · Datos: armonización a CIUO-08 en seis países}
Para comparar seis países hay que hablar el mismo idioma estadístico, y esa es la
parte menos glamorosa pero más laboriosa del proyecto.

Cada país tiene su propia encuesta de empleo, con su propia frecuencia y su propia
clasificación de ocupaciones. Costa Rica usa la ECE, trimestral. Colombia, la
GEIH, mensual, que agrego a trimestres. Ecuador, la ENEMDU, también mensual.
México, la ENOE, trimestral, donde aprovecho una correspondencia oficial del INEGI.
Uruguay, la ECH, con un trabajo de campo continuo que separo por mes de entrevista.
Y el Perú —destacado en la tabla— usa la EPEN, que es trimestral pero solo cubre
Lima Metropolitana, y que tiene además un cambio de clasificador que les explico en
la siguiente diapositiva.

A todas las llevo a una misma clasificación internacional, la CIUO-08 a cuatro
dígitos, y a una ventana común de tiempo: del primer trimestre de 2021 al cuarto de
2025. Esa ventana es deliberadamente posterior a la COVID, para no confundir el
efecto de los LLMs con el de la pandemia.

Tiene un costo: el período ``pre'', antes del tratamiento, es corto, apenas siete
trimestres limpios. Eso le quita potencia a las pruebas de tendencias previas, y
por eso más adelante me apoyo en las cotas de Honest-DiD para ser honesto sobre esa
limitación.
\end{frame}

% ===== 11 =====
\begin{frame}[allowframebreaks]{11 · Perú: el cambio de clasificador coincide con el shock}
Esta diapositiva trata el principal reto técnico del caso peruano, y quiero
explicarlo bien porque es delicado.

La EPEN cambió su clasificador de ocupaciones justo en el cuarto trimestre de 2022.
Hasta el tercer trimestre de 2022 usaba el clasificador CO-1995, a tres dígitos,
con 319 códigos. Desde el cuarto trimestre de 2022 pasó al CNO-2015, a cuatro
dígitos, con 473 códigos. El problema salta a la vista: el cambio de clasificador
cae exactamente en la misma fecha del tratamiento. Si no se trata con cuidado, uno
podría confundir un mero cambio de codificación con un efecto real de la tecnología.

¿Cómo lo resuelvo? Llevo ambos períodos a la clasificación internacional CIUO-08 a
cuatro dígitos, usando los cuatro anexos oficiales del INEI, pero con un refinamiento:
empleo pesos posteriores empíricos, condicionados en variables como la edad, el
sexo, la educación, el sector económico y el dominio geográfico. En palabras
simples: no traduzco a ciegas, sino tomando en cuenta el perfil del trabajador.

Y validé esa traducción con criterios fijados de antemano. El salto máximo en la
composición de ocupaciones entre el tercer y el cuarto trimestre de 2022 fue de
apenas dos puntos porcentuales, muy por debajo del umbral de cinco que me había
puesto. La divergencia máxima de los pesos entre 2023 y 2025 fue de cero coma cero
cuarenta y seis, también por debajo del umbral. Y queda pendiente, como robustez,
un DiD tipo \emph{donut} que descarte los trimestres del corte. En resumen: el
cambio de clasificador está bajo control, y puedo mostrar que no contamina el
resultado.
\end{frame}

% ===== 12 =====
\begin{frame}[allowframebreaks]{12 · Lima Metropolitana 2021: descriptivos}
Antes de los resultados, una foto de la muestra en el año base, 2021, para que
tengan en mente con quién estamos trabajando.

La edad promedio es de unos 38 años. Las mujeres son el 45 por ciento de la
muestra. Un 36 por ciento tiene educación terciaria o postsecundaria. La tasa de
formalidad, medida por afiliación al seguro de salud, es de un 37 por ciento. Se
trabajan, en promedio, unas 43 horas semanales, y el logaritmo del salario por hora
es de 1,88.

A la derecha está la distribución de la exposición $\beta$, ponderada por empleo.
La exposición media es de 0,28. Un dato que me parece relevante: una de cada cuatro
personas —el 25,6 por ciento— trabaja en ocupaciones con exposición cero, es decir,
prácticamente no tocadas por los LLMs.

Y hay un patrón que anticipa todo lo que viene: la formalidad sube de forma marcada
a medida que aumenta la exposición. En el quintil menos expuesto, la formalidad es
del 26 por ciento; en el más expuesto, sube al 51 por ciento. Dicho de otro modo,
en Lima las ocupaciones más expuestas a la IA son justamente las más formales. Por
eso el análisis por formalidad, que viene después, es tan central.
\end{frame}

% ===== 13 =====
\begin{frame}[allowframebreaks]{13 · Resultado 1: estimaciones agregadas para Perú}
Llegamos al primer resultado: el efecto agregado para el Perú. Y aquí hay una
lección importante sobre cómo se promedian los efectos.

En el Panel A está la especificación principal, el TWFE continuo, es decir, el
estimador de efectos fijos en dos vías. ¿Qué muestra? Sobre el empleo, un
coeficiente de 0,221, pero no significativo. Sobre el salario, 0,042, tampoco
significativo. Y sobre las horas, menos 3,25, marginalmente significativo. Si uno
se quedara solo con este panel, concluiría que ``no pasa nada'': sería la
replicación del nulo, como en Dinamarca.

Pero veamos el Panel B, que usa el estimador binario de Callaway y Sant'Anna,
partiendo la muestra en la mediana de exposición. Este estimador agrega
correctamente los efectos por cohorte y por tiempo. Y aquí aparece otra cosa: un
efecto positivo y muy significativo sobre el empleo, de 0,144, significativo al uno
por ciento; el salario sigue plano; y un efecto negativo y significativo sobre las
horas, de menos 0,8.

La moraleja es esta: el ``nada pasa'' del Panel A y el ``sí pasa'' del Panel B no
se contradicen por casualidad. Cuando uno agrega bien los efectos, aparece una
señal —más empleo, menos horas— que el promedio ingenuo del Panel A escondía. En la
siguiente diapositiva les explico por qué, y por qué prefiero el Panel B.
\end{frame}

% ===== 14 =====
\begin{frame}[allowframebreaks]{14 · Robustez 1: cota Honest-DiD}
Esta diapositiva es un ejercicio de honestidad metodológica, y quiero subrayarlo
porque habla de la seriedad del diseño.

Me impuse, de antemano, un estándar de robustez: para tomarme en serio una
afirmación principal, el resultado debía resistir cierta violación de las
tendencias paralelas. Técnicamente, usando el método de Rambachan y Roth, exigía un
valor de quiebre de al menos dos. En cristiano: el efecto tenía que sobrevivir
incluso si las tendencias previas no eran perfectamente paralelas.

¿Qué pasó con la especificación continua del Panel A? Que ese valor de quiebre fue
cero, para los tres resultados. Es decir, el intervalo de confianza original ya
incluía el cero, así que el criterio ni siquiera logra separar la especificación
continua de la hipótesis nula.

Lo digo con todas sus letras: la especificación continua TWFE no supera el estándar
de robustez que el propio artículo se impuso. Lo reporto de manera transparente.
Y por eso, a la luz de esto, trato al estimador binario de Callaway y Sant'Anna,
el del Panel B, como la lectura preferida del agregado.

Por último, esta discrepancia entre paneles no es un fracaso: es informativa. Cuando
hay respuestas heterogéneas a la dosis de tratamiento, el TWFE puede asignar pesos
implícitos negativos a algunos grupos y, al promediar, diluir o cancelar el efecto.
Es un fenómeno bien documentado, por de Chaisemartin y D'Haultfœuille, y por Sun y
Abraham. En otras palabras: el promedio simple miente, y por eso necesito un
estimador que agregue bien.
\end{frame}

% ===== 15 =====
\begin{frame}[allowframebreaks]{15 · Resultado 2: tendencias previas y dinámica del efecto}
Pasemos a las tendencias previas y a la dinámica del efecto en el tiempo, que es lo
que muestra el \emph{event study} de la izquierda.

Lo primero que hay que comprobar en un diseño de diferencias en diferencias es que,
antes del tratamiento, los grupos venían evolucionando en paralelo. El test
conjunto de Wald sobre los coeficientes previos no rechaza esa hipótesis en ninguno
de los tres resultados: los \emph{p}-valores están todos por encima de 0,17. Es
decir, no encuentro evidencia de tendencias previas problemáticas. Con la salvedad,
que ya mencioné, de que el período pre es corto y por tanto la prueba tiene poca
potencia.

Sobre el empleo, en el panel (a): los doce coeficientes posteriores al tratamiento
son positivos sin excepción, y van de más 0,25 a más 0,52. Es un patrón muy
consistente: tras la llegada de los LLMs, el empleo en las ocupaciones expuestas se
expande, y se mantiene expandido.

Sobre las horas, en el panel (c): los coeficientes previos son sistemáticamente
positivos, y los posteriores, negativos. Aquí debo ser cauto. Ese perfil de
``inversión'' —positivo antes, negativo después— sugiere una posible reversión de
tendencia que podría estar inflando mecánicamente el coeficiente principal de las
horas. Siguiendo la recomendación de Roth, lo tomo con pinzas y lo señalo como una
debilidad, no como un hallazgo firme.
\end{frame}

% ===== 16 =====
\begin{frame}[allowframebreaks]{16 · Resultado 3: heterogeneidad por formalidad}
Y llegamos al resultado que, para mí, es el más interesante de toda la charla: qué
pasa cuando separo el efecto por la formalidad de base de la ocupación.

Miren la tabla. En el canal formal, el efecto sobre el empleo es positivo y muy
significativo, 0,567; sobre el salario por hora es negativo y significativo, menos
0,128; y sobre las horas, negativo. En el canal informal, en cambio, el cuadro se
invierte: el empleo no se mueve de forma significativa, pero el salario por hora
sube de forma fuerte y significativa, 0,265.

Es decir, el sector formal y el informal se mueven en direcciones opuestas. Y esas
diferencias no son ruido: la prueba de igualdad entre ambos canales se rechaza para
el empleo, con un \emph{p}-valor de 0,005, y para el salario, con un \emph{p}-valor
menor a 0,001.

La conclusión es contundente, y reordena todo lo anterior: aquel ``nulo agregado''
que veíamos en el Panel A es, en realidad, un artefacto de cancelación. No es que no
pase nada; es que pasan dos cosas opuestas —una en el formal y otra en el informal—
que, al promediarse a ciegas, se anulan entre sí. Solo al separar por formalidad
aparece la verdadera estructura del efecto.
\end{frame}

% ===== 17 =====
\begin{frame}[allowframebreaks]{17 · Heterogeneidad por formalidad: visualización}
Esta diapositiva es simplemente la versión visual del resultado anterior, para que
quede grabado de un vistazo.

En el gráfico ven los coeficientes del canal formal y del canal informal, cada uno
con su intervalo de confianza al noventa y cinco por ciento. Lo que quiero que noten
es la separación: para el empleo y para el salario por hora, los dos canales no solo
apuntan en direcciones distintas, sino que sus intervalos casi no se traslapan. Esa
separación visual es la misma que la tabla resumía con los \emph{p}-valores de
0,005 para el empleo y menor a 0,001 para el salario.

En una sola frase: la formalidad no es un detalle; es el eje sobre el que se
organiza todo el efecto de los LLMs en el caso peruano.
\end{frame}

% ===== 18 =====
\begin{frame}[allowframebreaks]{18 · Síntesis: el ``Patrón Peruano'' emergente}
Recojo todo lo anterior en una síntesis, y empiezo siendo honesto sobre las cuatro
narrativas que había precomprometido: ninguna de las cuatro reproduce exactamente lo
que encontré en Lima.

La narrativa A requería que el efecto formal sobre el empleo fuera negativo, y me
salió positivo. La B requería que ambos canales fueran nulos, y rechazo esa igualdad
tanto en empleo como en salario. La C requería que el canal formal fuera
aproximadamente nulo, pero la magnitud y la significancia del canal formal lo
contradicen. Y la D predecía ganancias salariales para los formales, cuando el signo
salarial formal me salió, de hecho, opuesto.

Entonces, ¿qué encontré? Lo que llamo, de manera exploratoria, el ``Patrón
Peruano''. Las ocupaciones formales-dominantes expuestas a los LLMs expanden el
empleo y reducen las horas, con una presión a la baja sobre el salario por hora. Las
ocupaciones informales-dominantes, en cambio, muestran una contracción débil del
empleo, pero con salarios al alza entre quienes permanecen en la ocupación, lo cual
podría reflejar selección positiva o un efecto de equilibrio general.

Como este patrón no estaba entre los cuatro pre-registrados, lo propongo como un
quinto patrón, que añadiré a la pre-registración formal antes de mirar los datos
posteriores de los otros cinco países. Así mantengo la disciplina del diseño.
\end{frame}

% ===== 19 =====
\begin{frame}[allowframebreaks]{19 · Puntos a considerar (I): identificación y datos}
Quiero dedicar dos diapositivas a las limitaciones, con toda transparencia, porque
creo que reconocerlas es parte de hacer buena investigación. Empiezo por las de
identificación y datos.

Primero, la cobertura geográfica. La EPEN cubre únicamente Lima Metropolitana, así
que mi estimación para el Perú no es nacionalmente representativa. Tengo prevista una
validación con la ENAHO anual, que sí es nacional.

Segundo, la definición de formalidad. La construyo a partir de la afiliación al
seguro de salud. La definición de la OIT, basada en pensiones, requeriría también la
ENAHO.

Tercero —y ya lo discutí— la cota Honest-DiD es cero en el TWFE continuo. La
especificación principal no supera el criterio de robustez que me había puesto, y por
eso la lectura preferida pasa al estimador binario de Callaway y Sant'Anna.

Y cuarto, la posible reversión de tendencia en las horas, que ya señalé en el
\emph{event study}: los coeficientes previos positivos podrían estar inflando el
coeficiente de horas, así que lo trato con cautela.
\end{frame}

% ===== 20 =====
\begin{frame}[allowframebreaks]{20 · Puntos a considerar (II): supuestos y método}
Sigo con las limitaciones de supuestos y método.

Quinto, la selección no modelada. La EPEN es un corte transversal repetido, no un
panel. Eso significa que la migración entre ocupaciones, la salida hacia el
desempleo y la recomposición demográfica no son identificables con estos datos. Veo
fotos sucesivas de la población, no sigo a las mismas personas.

Sexto, el \emph{crosswalk} descansa en pesos posteriores empíricos. Queda pendiente
la escalera completa de robustez: probar con CO-1995 a tres dígitos, el DiD
\emph{donut}, pesos uniformes y una versión solo con datos posteriores.

Séptimo, el supuesto SUTVA. Si hay derrames dentro de las empresas, mi estimación
del desplazamiento estaría sesgada hacia cero. Una cota narrativa, de Berg y
coautores, acota la magnitud de ese problema, pero no lo elimina.

Octavo, el período previo es corto: solo siete trimestres. Las pruebas de tendencias
previas tienen, por tanto, poca potencia.

Y noveno, lo repito una vez más: las estimaciones de Perú son previas al depósito del
pre-registro, así que las identifico como exploratorias. El pre-registro formal en
OSF cubrirá los cinco países restantes.
\end{frame}

% ===== 21 =====
\begin{frame}[allowframebreaks]{21 · Próximos pasos}
¿Hacia dónde va el proyecto a partir de aquí?

En lo inmediato, voy a cerrar la escalera de robustez del clasificador para el Perú,
y a validar los resultados con la ENAHO anual, que me dará representatividad nacional
y la definición de formalidad de la OIT. También probaré medidas alternativas de
exposición —además de la $\beta$ de Eloundou, las medidas de Webb y el índice AIOE de
Felten y coautores— y añadiré las cotas de Honest-DiD, el método de Oster y los
errores estándar adecuados para el enfoque \emph{shift-share}.

Y en el horizonte central del proyecto: replicar todo este \emph{pipeline} en Costa
Rica, Colombia, Ecuador, México y Uruguay. La meta final es poner los seis países
lado a lado y comparar sus efectos. Esa comparación entre países es, como dije al
inicio, la contribución empírica central del artículo. Lo de hoy, el caso peruano,
es la primera pieza de ese rompecabezas.
\end{frame}

% ===== 22 =====
\begin{frame}[allowframebreaks]{22 · Algunas referencias clave}
Esta última diapositiva recoge las referencias clave del trabajo, que ya fui citando
a lo largo de la charla: el marco de tareas de Acemoglu y Restrepo; las medidas de
exposición de Eloundou y coautores; los estimadores de diferencias en diferencias de
Callaway, Goodman-Bacon y Sant'Anna; las herramientas de robustez de Rambachan y Roth
y de Roth; y los antecedentes sobre América Latina, incluido el trabajo del BID del
que soy coautor.

Permítanme cerrar con una sola idea. La pregunta de si la inteligencia artificial
generativa está cambiando el trabajo no se responde igual en todas partes. En una
región con altísima informalidad y mercados segmentados como la nuestra, el efecto
puede tener una estructura propia. La evidencia preliminar para el Perú sugiere,
justamente, que el promedio agregado esconde dos fuerzas opuestas según la
formalidad. Confirmar si ese patrón se repite, o cambia, en los otros cinco países
es el siguiente paso.

Muchas gracias por su atención. Quedo atento a sus comentarios y sugerencias.
\end{frame}

\end{document}
